%---------- 第一章 绪论 ----------
\section{绪\quad 论}

\subsection{研究背景与意义}

\quad 图像识别是计算机视觉领域的核心任务之一，广泛应用于人脸识别、医学影像分析、
工业质检、自动驾驶等场景。传统图像识别方法主要依赖人工设计的特征（如 SIFT、HOG）
与浅层分类器，特征表达能力有限，在复杂场景下的识别效果难以满足实际需求。

近年来，以卷积神经网络（Convolutional Neural Network, CNN）为代表的深度学习方法
通过端到端的特征自动学习，大幅提升了图像识别的精度，成为该领域的主流技术
路线\cite{ref-lecun2015}。研究基于深度学习的图像识别方法，对推动相关应用落地具有
重要的理论意义与实用价值。

\subsection{国内外研究现状}

\quad 自 2012 年 AlexNet 在 ImageNet 图像分类竞赛中取得突破以来，VGGNet、GoogLeNet、
ResNet\cite{ref-he2016} 等一系列深层网络结构相继被提出，网络深度与识别精度不断提升。
国内研究者在中文场景识别、轻量化网络设计等方向也取得了丰富成果\cite{ref-wang2020}。
然而，现有方法在目标尺度变化剧烈、背景干扰复杂的场景下仍存在识别精度不足的问题。

\subsection{本文主要工作}

\quad 本文的主要工作包括：
\begin{enumerate}
  \item 梳理卷积神经网络的基本原理与典型结构，分析现有图像识别方法的优势与不足；
  \item 提出一种结合注意力机制的改进网络模型，增强网络对关键区域的特征提取能力；
  \item 在公开数据集上开展对比实验，验证所提方法的有效性。
\end{enumerate}

\subsection{论文组织结构}

\quad 本文共分为四章：第一章为绪论，介绍研究背景、意义与国内外研究现状；第二章介绍
相关理论与技术基础；第三章详细阐述所提方法并给出实验结果与分析；第四章为总结与展望。
